% Section 5 (writer: steep-b). Source: research/tracks/steeper/notes-final.md §2.5-§2.6 (Props 2.9-2.17, Lemma 2.11a,
% end of §2.6), §4.4. Notation: ../paper/NOTATION.md; St's guard noise eta -> epsilon_G, step count M -> n_st,
% input law pi -> mu, slack eta -> epsilon. The test set of Prop 2.16 is written Z (S and E are the error rule and set).
\section{Limits of the trick}\label{sec:limits}

\Cref{thm:teacher:main} needs $\lambda$ to grow linearly in $n$ and the error rule to be incompressible relative to the model. Neither requirement can be dropped, and steeper selection by description length has further failure modes. The setting is \cref{def:teacher:instance} unless stated otherwise; proofs are in \cref{app:limits}.

\subsection{Why \texorpdfstring{$\lambda$}{lambda} must be proportional to \texorpdfstring{$n$}{n}}\label{sec:limits:schedule}

\begin{proposition}[Plain Bayes, MDL and SI learn the errors]\status{proved; computed}\src{steeper Props 2.9, 2.10; checks c3, c6(G), c7(G)}\label{prop:limits:plain}
\begin{enumerate}[label=(\alph*)]
\item Let $P$ be a shared hypothesis with $P\ge2^{-\ell(\pi)-\kappa}\nu_\pi$ for every $\pi\in\Hcl{\times}$, such as $\Mprod$ or $\MF$. For every sequence of distinct inputs and every $n$, $\sum_{i\le n}-\log P(y^*(x_i)\mid x_i,D_{i-1})\le\ell(\pi_b)+\kappa\le d+\lceil L_S\rceil+\kappa_b+\kappa$. So at most $d+L_S+O(\log dm)$ inputs in total get probability below $\frac12$ for the teacher's label (semimeasure convention, hence also normalised), and every other revealed error input gets the erroneous label with probability at least $\frac12$; if all inputs are revealed, this is a fraction $1-9.6\cdot10^{-24}$ of the error inputs at H\"anni's sizes ($\kappa=0$).
\item Let $B<\infty$, let $\lambda=\lambda_n\ge1$ be fixed or depend on $n$, and draw $n\le2^d$ inputs uniformly without replacement. With probability at least $1-\delta$, every $\lambda$-MAP $\hat\pi$ has $R_B(\hat\pi)\le\lambda\hat L/n+B\sqrt{(\hat L\ln2+\ln(2/\delta))/(2n)}$ with $\hat L:=d+\lceil L_S\rceil+\kappa_b$, hence $e(\hat\pi)\ge2^{-R_B(\hat\pi)/p}-2^{-B}$. If $\lambda_n/n\to0$, this tends to $1-2^{-B}$ as $n$ grows (within $n\le2^d$).
\item On the interleaved string with i.i.d.\ uniform inputs, Solomonoff's $M$ has expected total label loss at most $d+L_S+O(\log dm)$, and the steeper SI predictor pays at most $\lambda\ell_U$ more (\cref{thm:absorb:absorption}(d)).
\end{enumerate}
\end{proposition}

\emph{Proof idea.} (a) Telescope $P(D_n)\ge2^{-\ell(\pi_b)-\kappa}$; for $\MF$ this is the mistake bound of \citet{claudeSFI2026fundamentals}. (b) is \cref{lem:teacher:finite} with the memoriser as reference. At $\kappa=0$, $B=30$, $\delta=0.01$, a direct computation (code in the accompanying repository) gives $e(\hat\pi)\ge0.649$ at $n=10^{13}$ for $\lambda=1$, and $e(\hat\pi)\ge0.734$ at $n=10^{15}$ for $\lambda_n=\sqrt n$, whatever $\delta_S$ (\cref{app:limits:plain}).

\begin{remark}[Fixed $\lambda$ versus $\lambda\propto n$]\status{remark (from proved results); known}\src{steeper end of \S2.5, \S4.4}\label{rem:limits:schedule}
So the trick needs $\lambda=\Theta(n)$: $\lambda/n$ must stay above $2^{-m}$ (\cref{thm:teacher:main}(iii),(iv)), and every schedule with $\lambda_n/n\to0$ eventually reproduces the errors. This confirms H\"anni's remark that ``we will need to increase the hypothesis term coefficient compared to the sum of log likelihoods term as the data set size increases'' \citep{hanni2026lesswrong}. A fixed $\lambda>1$ removes noise (\cref{thm:split:private}), not finite-complexity errors. With $\lambda=cn$ the learner discards every component whose rate is below $c$, valid or not (\cref{prop:limits:separable}): it is deliberately inconsistent for the teacher's label distribution. Its price is a batch excess: $-\log\xi_\lambda(D_n)\le n\,(c\,\ell(\pi)+\hat R(\pi))$ for every $\pi$, with $\hat R$ the empirical risk, so the code length exceeds $\pi$'s own loss by $c\,\ell(\pi)$ bits per sample, against $\ell(\pi)/n$ at $\lambda=1$. This is not a sequential regret, since $\lambda=cn$ is not one Bayesian predictor (\cref{rem:setup:persample}). That a bounded $\lambda$ overfits and $\lambda$ must grow with the sample size is known for $\lambda$-weighted two-part MDL in noisy classification \citep{zhu2025quantifying,zhu2026overfitting} and is not claimed here (\cref{rem:related:mdl}); there, consistency means converging to the data distribution, which here contains the teacher's errors, and by (b) every schedule with $\lambda_n/n\to0$ learns them.
\end{remark}

\subsection{Which error rules defeat the trick}\label{sec:limits:rules}

H\"anni's example fixes only a density and a description length, and realistic systematic errors are often low-complexity patterns. Split the rule as $S=S_1\sqcup S_2$, $|S|=s_0=p2^m$, with $s_i:=|S_i|$, $p_1:=s_1/2^m$ and $p_2:=s_2/(2^m-s_1)$, where $S_1$ is \emph{cheap}, $K(S_1\mid m)\le k_1$; every input of $E_1:=S_1\times\{0,1\}^{d-m}$ is an error input. The \emph{envelope} $\pi_e$ predicts $f_a\oplus1$ with probability $1$ on $E_1$ and $f_a$ with flip probability $p_2$ elsewhere; it has risk $R_e=(1-p_1)H(p_2)<H(p)$ and length at most $d+k_1+\kappa_e$, $\kappa_e:=2\log d+2\log m+2\log(s_2+1)+\kappa$. Put $J_e(c):=c(d+k_1+\kappa_e)+R_e$, $\zeta:=c(k_1+\kappa_e+4+\Lambda_1)+\epsilon$ and $c_{\rm env}:=(H(p)-R_e)/(k_1+\kappa_e+4+\Lambda_1)$.

\begin{proposition}[Error envelopes]\status{proved; computed}\src{steeper Prop 2.17 (T5 Prop 4.5, all-hypothesis form); check c7(B); referee q4(C)}\label{prop:limits:envelope}
Let $B=\infty$, and let $J_c(\pi)\le J_e(c)+\epsilon$ with $J_e(c)+\epsilon<\Phi_2(c)$.
\begin{enumerate}[label=(\roman*)]
\item $\lev(\pi)=1$ and $R(\pi)\le R_e+\zeta$. So for $c<c_{\rm env}$ and $\epsilon$ small enough that $\zeta<H(p)-R_e$, no program with $R(\pi)\ge H(p)$, the good model included, is $\epsilon$-near-optimal.
\item With $B_2:=K(S_2\mid a,K(a))$, the average probability of reproducing the cheap errors is
\begin{gather*}
e_1(\pi):=\Ex\big[\nu_\pi(y^*(x)\mid x)\,\big|\,x\in E_1\big]\ \ge\ 2^{-\bar R_1},\\
\bar R_1:=\tfrac1{p_1}\big(\zeta+(1-p_1)H(p_2)-\min(c,2^{-m})\,(B_2-K(S_1))^+\big).
\end{gather*}
\end{enumerate}
\end{proposition}

\emph{Proof idea.} (i) is \cref{thm:teacher:main}(o) with (D1). For (ii), code $S_2$ given $S_1$ with the hypothesis, as in \cref{lem:teacher:coding}; this bounds its risk off $E_1$ from below, so its risk on $E_1$ is small (\cref{app:limits:envelope}).

At H\"anni's sizes take $S_1=\{z:z_1=\dots=z_j=0\}$ and $S_2$ uniformly random. For $j=13$ ($128$ cheap and $896$ random points; $\log\binom{2^m-s_1}{s_2}=10{,}418$ bits), a direct computation gives $c_{\rm env}=1.64\cdot10^{-5}$, which is $31\%$ of the logarithmic width of the window $(9.54\cdot10^{-7},8.06\cdot10^{-3}]$ of \cref{thm:teacher:main} at $\kappa=0$ ($25\%$ at $\kappa=30$), and $e_1\ge0.545$ at $c=1.01\cdot2^{-m}$ ($\kappa=0$, $B_2$ typical). For $j=11$ ($512+512$ points, $6{,}364$ bits): $47\%$ of the log-width and $e_1\ge0.862$ (table in \cref{app:limits:envelope}).

So \emph{the extension of H\"anni's example to every error rule with his density and description length is false}: the rule with $j=13$ has density $2^{-10}$ and about $10{,}400$ bits, close to his ``$10000$ bits'', yet on almost a third of the window's log-width no hypothesis that fits no better than the good model is near-optimal, and near the lower end near-optimal hypotheses reproduce the cheap errors. (Its deficiency is about $1{,}300$ bits, so \cref{thm:teacher:main} does not cover it; above $c_{\rm env}$ nothing is proved for it.) H\"anni called the example ``definitely somewhat cherry-picked for the idea to work, though I also don't consider it completely contrived'' \citep{hanni2026lesswrong}. What is confirmed is the example for error rules incompressible relative to the model. For structured rules the dangerous hypotheses are envelopes, not universal ones: a cheap, dense error region has a high rate and is learned like any valid component. The extreme case is a cheap, frequent rule.

\begin{proposition}[Rate blindness]\status{proved; computed}\src{steeper Prop 2.11; checks c3, c6(A)}\label{prop:limits:rateblind}
Let $B=\infty$ and $\ell(\pi_{\rm coin})\le\kappa$.
(i) If $H(p)/(\ell(\pi_b)-\ell(\pi_g))\ge(1-H(p))/(\ell(\pi_g)-\ell(\pi_{\rm coin}))$, then at every $c$ the good model is beaten by the memoriser or beaten or tied by the coin.
(ii) If
\[
\frac{H(p)}{\lceil L_S\rceil+\kappa_b+4+\Lambda_1}\ \ge\ \frac{1-H(p)}{d-4-\Lambda_1-\kappa}\quad\text{and}\quad d-4-\Lambda_1-\kappa>0,\tag{$\star$}
\]
then at every $c$ every program with $u>\frac34$ and $R\ge H(p)$ is beaten by the memoriser or beaten or tied by the coin.
(iii) For $c<2^{-m}$, near-optimal programs reproduce the errors, under the condition of \cref{thm:teacher:main}(iii).
\end{proposition}

\begin{lemma}[The criterion is robust in the machine constant]\status{proved; computed}\src{steeper Lemma 2.11a; check c6(A)}\label{lem:limits:robust}
Fix $m$ and $r$, let $\kappa$ bound the machine constants in $\Lambda_1$, $\kappa_b$ and $\ell(\pi_{\rm coin})$, and put $H:=H(2^{-r})$, $q:=(1-H)/H$, $a:=2+4q$, $b:=4+2\log m+2\kappa+q(\lceil L_S\rceil+4\log m+2\log s_0+2\kappa+4)$ and $F(d):=d-a\log d-b$. (i) For $d\ge1$, $(\star)$ holds if and only if $F(d)\ge0$. (ii) If $b\ge a/\ln2$, then $\{d\ge1:F(d)\ge0\}=[d_0,\infty)$ for a unique real $d_0=d_0(\kappa)$. (iii) $d_0(\kappa)$ grows linearly in $\kappa$, with slope $(2+2q)(1+o(1))$.
\end{lemma}

For a single $3$-bit prefix as error rule ($m=3$, $p=1/8$), a direct computation gives $d_0=49$ at $\kappa=0$ and $169$ at $\kappa=30$; for a single $6$-bit prefix, $450$ and $1005$ (\cref{app:limits:robust}). So for each fixed cheap rule and each $\kappa$, once the simple model is long enough, no program that knows the parity and fits no better than noise is optimal at any $c$. H\"anni's $11{,}710$-bit rule, by contrast, leaves a window of a factor of about $8000$ ($\kappa=0$) or $3800$ ($\kappa=30$).

\subsection{Other pathologies}\label{sec:limits:other}

\begin{proposition}[Separable components]\status{proved; computed}\src{steeper Prop 2.12 (T5 Thm 3.2, Thm 4.1); check c5(a)}\label{prop:limits:separable}
Let the inputs split into regions $X_1,\dots,X_k$ and the hypotheses be products of per-region choices with additive code lengths. Each region selects a vertex of its own hull for slope $c$ (\cref{lem:setup:hull}). If region $j$ offers a default or a component $j$ of extra length $k_j$ and population risk reduction $r_j$, component $j$ is selected if and only if $r_j/k_j>c$ (ties arbitrary). So a valid component whose rate is at most an error component's rate cannot be kept without the error, except at a common tie value. Two worlds with the same teacher channel and input law, one where a component is an error and one where it is a rare valid rule, get the same selection at every $c$.
\end{proposition}

The steeper penalty sorts components by \emph{rate}, risk reduction per bit, not by validity. Three consequences follow (proofs in \cref{app:limits:guard}).

\begin{proposition}[Guard erosion]\status{proved; computed}\src{steeper Prop 2.13 (T5 Prop 4.4); checks c5(b), c6(G), c7(E); referee q1(I)}\label{prop:limits:guard}
In a region of contexts a schema $\sigma$ is one of $n_{\rm st}$ possible steps, and a guard $G$ splits the region into $G$-contexts (input mass $\mu(G)$), where the teacher plays $\sigma$ with probability $1-\epsilon_G+\epsilon_G/n_{\rm st}$ and every other step with probability $\epsilon_G/n_{\rm st}$, and $\neg G$-contexts (mass $\mu(\neg G)$), where it plays one fixed step other than $\sigma$. Over a common base, ``none'' is uniform (length $0$), ``unguarded $\sigma$'' predicts the $G$-context law everywhere (length $k_\sigma$), and ``guarded $\sigma$'' does so in $G$-contexts and is uniform over the other $n_{\rm st}-1$ steps elsewhere (length $k_\sigma+k_G$). With $H_G$ the entropy of the $G$-context law, $g:=\log n_{\rm st}-H_G$ and $\Delta_G:=\log(n_{\rm st}/\epsilon_G)-\log(n_{\rm st}-1)$, put $\rho_\sigma:=(\mu(G)g-\mu(\neg G)\log(1/\epsilon_G))/k_\sigma$ and $\rho_G:=\mu(\neg G)\Delta_G/k_G$. If $\rho_G<\rho_\sigma$, then for every $c\in(\rho_G,\rho_\sigma)$ the unguarded rule is the unique selection, and it applies $\sigma$ where the teacher never did. Otherwise the selection goes from guarded to none.
\end{proposition}

With $n_{\rm st}=64$, $\epsilon_G=0.05$, $\mu(G)=0.05$, $\mu(\neg G)=5\cdot10^{-4}$, $k_\sigma=22$ and $k_G=12$, a direct computation gives $\rho_\sigma=1.2226\cdot10^{-2}$ and $\rho_G=1.8103\cdot10^{-4}$: the unguarded rule wins on a factor of $67$ in $c$.

\begin{proposition}[$\lambda$ cannot be chosen from imitation data]\status{proved; computed; (i) known}\src{steeper Prop 2.14 (T5 Thm 4.3); check c5(c)}\label{prop:limits:noident}
(i) Every strictly proper scoring rule has expected score on the teacher's data uniquely maximised by the teacher's channel $P^*(\cdot\mid x)$, errors included \citep{gneiting2007strictly}, so validation by such a score prefers the teacher's channel to the intended model. (ii) Along the population selection path the held-out log loss $R(h_c)$ is nonincreasing as $c$ decreases (\cref{lem:setup:hull}(ii)), so log-loss validation over a grid of values of $c$ is minimised at the smallest, error-memorising, value; in H\"anni's example the losses are $1.0$, $0.0112$, $0.0112$, $0$ at $c=10^{-1},10^{-3},10^{-5},10^{-7}$. (iii) In the twin worlds of \cref{prop:limits:separable}, every rule that chooses $c$ from imitation data alone is wrong in one of them.
\end{proposition}

\begin{proposition}[The selection depends on the input law]\status{proved}\src{steeper Prop 2.15}\label{prop:limits:inputlaw}
In the setting of \cref{prop:limits:separable}, let the inputs have law $\mu$ and let $\Delta_j$ be component $j$'s expected log-loss reduction per input of $X_j$. Then component $j$ is selected if and only if $\mu(X_j)\Delta_j>c\,k_j$. So the window that keeps a valid component $j$ and drops an error component $j'$ is $c\in(\mu(X_{j'})\Delta_{j'}/k_{j'},\,\mu(X_j)\Delta_j/k_j)$; it moves with $\mu$ and is empty once $\mu(X_{j'})/\mu(X_j)\ge\Delta_jk_{j'}/(\Delta_{j'}k_j)$, and after a shift to $\mu'$ a $c$ chosen in $\mu$'s window keeps the error if and only if $\mu'(X_{j'})\Delta_{j'}>c\,k_{j'}$.
\end{proposition}

\Cref{thm:teacher:main}'s window depends on the input law by the same mechanism; this is proved here only in the separable case.

\subsection{Unseen properties}\label{sec:limits:ood}

For $\MF$, a simple property absent from the training data gets a constant probability of bad behaviour, $2^{-K(Q,b)-O(1)}$, where $K(Q,b)$ is the length of a shortest self-delimiting code for a decider of a set $Q$ and a program for a behaviour $b$ \citep[H\"anni's proposition, with no log term]{claudeSFI2026fundamentals}. The steeper counterpart is as follows (proof in \cref{app:limits:ood}).

\begin{proposition}[Steeper simplicity and unseen properties]\status{proved}\src{steeper Prop 2.16}\label{prop:limits:ood}
In $\Hcl{\times}$, let $Q\subseteq X$ be decidable, $b$ partial computable, $D$ a data set with $\dom D\cap Q=\emptyset$, and $Z\subseteq Q\cap\dom b$ finite. For every $\lambda\ge1$: (i) the programs that answer $b(x)$ on $Q$ and run another program elsewhere carry posterior mass exactly $2^{-\lambda(K(Q,b)+\kappa)}$, so $\xi_\lambda(D\cup b|_Z)\ge2^{-\lambda(K(Q,b)+\kappa)}\xi_\lambda(D)$; (ii) if $\nu_{\pi_1}(D)=\nu_{\pi_2}(D)>0$, then $\Pi_\lambda(\pi_1\mid D)/\Pi_\lambda(\pi_2\mid D)=2^{-\lambda(\ell(\pi_1)-\ell(\pi_2))}$.
\end{proposition}

At $\lambda=1$, (i) is H\"anni's proposition for $\Mprod$. With $\lambda=cn$ the bound vanishes, so steeper simplicity removes his constant lower bound, but only through (ii): the posterior concentrates on the shortest hypotheses that fit the training data, whatever they do on $Q$. If the shortest one extrapolates well, steeper SFI is more reliable out of distribution than SFI; if a shorter one fits the training data and is wrong on $Q$, it is more confidently wrong. No guarantee about unseen properties follows.

\begin{remark}[A limitation of the method]\status{remark; computed}\src{steeper end of \S2.6; check c7(C)}\label{rem:limits:method}
The window needs $G_1>H(p)$, that is, $p<0.056$. The sharper floor of \cref{lem:teacher:sharpfloor} gives the kink for every $p<1/4$ (\cref{prop:algstat:parity}), but does not extend the window as the argument stands: with it and no deviation term the level-$2$ floor $1-\log(3/2-p)$ exceeds $H(p)$ only for $p<0.123$; and a program of level at least $2$ has $I(\pi)$ bounded only through its length, which near $c=2^{-m}$ allows $I(\pi)\approx L_S$ and a deviation term of order $\sqrt{H(p)/p}$. A window for $p<0.123$ and $c$ well above $2^{-m}$ looks within reach but is not proved.
\end{remark}

\paragraph{Open question.}\label{q:limits:open}
Which structured error rules leave a window? For envelopes with error density $\theta<1$ in the cheap region only part (i) of \cref{prop:limits:envelope} carries over; whether near-optimal hypotheses memorise the random part $S_2$, and what is selected above $c_{\rm env}$, are open. A theory for realistic, low-complexity human errors is the main gap between the trick and H\"anni's motivation. Is there a principled, data-free choice of $c$? By \cref{prop:limits:separable,prop:limits:noident} imitation data alone cannot tell errors from rare valid rules; what extra information would suffice is open.
